% BANKS-SOURCE: pdf=40 print=30 kind=prose
1PI expansion of connected Green functions then tells us that every $W_n$ has a pole
on each of its external legs. We will see that the residue of this multiple pole, after
proper normalization, contains the S-matrix for all scattering processes in which the
total number of ingoing and outgoing particles is $n$.

Before concluding this subsection, we want to remind the reader that our notation has
conflated the quantum field operator $\phi$, the functional integration variable $\phi$, and the
argument of the quantum action $\phi$. These are distinct concepts, but it is typographically
insane to use different versions of the same Greek letter to separate them. The reader
who has truly understood the discussions above should have no trouble distinguishing
the meaning of $\phi$ from its context.

% BANKS-SOURCE: pdf=40 print=30 kind=section id=3.6
\subsection{The Källen--Lehmann spectral representation}

% BANKS-SOURCE: pdf=40 print=30 kind=prose
For $x^0>0$, we can write the two-point function of the interacting Heisenberg field
$\phi(x)$ as

% BANKS-SOURCE: pdf=40 print=30 kind=equation id=3.39
\begin{equation}
\langle0|T\phi(x)\phi(0)|0\rangle
=\int\dd^D p\,\sum_n\delta^D(p-p_n)\ee^{+\ii p\mathbin{\cdot}x}\abs{\langle0|\phi(0)|n\rangle}^2.
\label{eq:3.39}
\end{equation}

We have inserted a complete set of states between the two fields, and used the action of
the translation generator on the Heisenberg fields and on the states. $p_n$ is the momentum
of the state $\ket{n}$. As we will discuss in more detail below, we assume that the theory has
a complete set of \,\emph{scattering states}, i.e. states that are composed of multiple stable
particles traveling freely at large space-like separation from each other. The energies
and momenta of such states are determined by the usual free-particle relations. We will

% BANKS-SOURCE: pdf=41 print=31 kind=prose
assume further that the theory has a mass gap. That is, apart from the non-degenerate
vacuum state, the lowest value of the invariant mass squared $-P^2$ is a non-zero number
equal to the mass squared of the lightest stable particle in the theory.

For theories with massless particles, the general form of the Källen--Lehmann \cite{banks-ref-14,banks-ref-15}
representation will remain unchanged, but the identification of particle masses with
poles is slightly more subtle, and depends in detail on the form of the massless particle
interactions. Let $m$ be the mass of the lightest stable particle that can be created from
the vacuum by a single action of $\phi$. Then it follows from Lorentz invariance that

% BANKS-SOURCE: pdf=41 print=31 kind=equation id=3.40
\begin{equation}
\langle0|\phi(0)|p\rangle=\sqrt{\frac{Z}{2\omega_{\mathbf p}(2\pi)^{D-1}}}.
\label{eq:3.40}
\end{equation}

The constant $Z$ is called the \emph{on-shell wave-function renormalization constant} of $\phi$.
In any field theory, with or without a mass gap, we will have only states of time-like
or null momentum, with positive energy. So we can rewrite the two-point function as

% BANKS-SOURCE: pdf=41 print=31 kind=equation id=3.41
\begin{equation}
\begin{split}
\langle0|T\phi(x)\phi(0)|0\rangle
={}&\int_0^\infty\dd\mu^2\int\dd^D p\,\theta(p^0)\delta(p^2+\mu^2)\\
&\quad\times\sum_n\delta^D(p-p_n)\ee^{+\ii p\mathbin{\cdot}x}\abs{\langle0|\phi(0)|n\rangle}^2.
\end{split}
\label{eq:3.41}
\end{equation}

We recognize that this is

% BANKS-SOURCE: pdf=41 print=31 kind=equation id=3.42
\begin{equation}
\langle0|T\phi(x)\phi(0)|0\rangle=\int_0^\infty\dd\mu^2\,\rho(\mu^2)D_F(x;\mu^2),
\label{eq:3.42}
\end{equation}

where $D_F$ is the free-field two-point function and

% BANKS-SOURCE: pdf=41 print=31 kind=equation id=3.43
\begin{equation}
\sum_n\delta^D(p-p_n)\abs{\langle0|\phi(0)|n\rangle}^2\equiv\rho(-p^2).
\label{eq:3.43}
\end{equation}

Indeed, since $U^\dagger(\Lambda)\phi(0)U(\Lambda)=\phi(0)$, for every Lorentz transformation
$\Lambda$, the positive function $\rho(-p^2)$ is Lorentz-invariant, and depends only on the
invariant mass squared $-p^2$.
In a theory with a mass gap, the lower limit on the integral is $m^2$, coming from the
one-particle state, and then there is a gap until the lightest multiparticle state that can
be created by the action of $\phi$. Thus there is a contribution

% BANKS-SOURCE: pdf=41 print=31 kind=display
\[
\rho(\mu^2)=Z\delta(\mu^2-m^2)+C(\mu^2),
\]

where the continuum contribution $C$ is positive and has support starting at the invariant
mass squared of the lowest multiparticle state.

It follows that the exact momentum-space two-point function $W_2(p)$ has a pole at
the stable particle mass, with residue $Z$. Note that nothing in this derivation requires
that the field $\phi$ be the fundamental variable that appears in the Lagrangian, or that
perturbation theory be applicable. Any field with a non-zero $Z$ will do. In cases where
perturbation theory is applicable, the fundamental fields will have poles in their Green
functions. However, in quantum chromodynamics (QCD), the theory of the strong
interactions, perturbation theory is inapplicable in the vicinity of hadron mass scales.
The fundamental quark and gluon fields are not associated with stable particles, and

% BANKS-SOURCE: pdf=42 print=32 kind=prose
hadrons are instead created by multi-linear functions of the quark fields. In lattice
gauge theory, the functional integral of QCD is done numerically, and hadron masses
are computed by finding poles in the two-point functions of quark--anti-quark bilinears
and quark trilinears.
